\documentclass[
    11pt,
    a4paper
]{article}

% =========================================================
% IDIOMA, CODIFICACIÓN Y TIPOGRAFÍA
% =========================================================

\usepackage[utf8]{inputenc}
\usepackage[T1]{fontenc}
\usepackage{csquotes}
\usepackage[french]{babel}
\usepackage{lmodern}
\usepackage{microtype}

% =========================================================
% GEOMETRÍA DE LA PÁGINA
% =========================================================

\usepackage[
    top=2.5cm,
    bottom=2.5cm,
    left=2.8cm,
    right=2.8cm
]{geometry}

% =========================================================
% MATEMÁTICAS
% =========================================================

\usepackage{amsmath}
\usepackage{amssymb}
\usepackage{amsfonts}
\usepackage{mathtools}
\usepackage{bm}

% =========================================================
% UNIDADES Y NOTACIÓN CIENTÍFICA
% =========================================================

\usepackage{siunitx}

\sisetup{
    output-decimal-marker={,},
    separate-uncertainty=true,
    per-mode=symbol
}

% =========================================================
% FIGURAS
% =========================================================

\usepackage{graphicx}
\usepackage{float}
\usepackage{subcaption}
\usepackage{placeins}

% Carpeta en la que se almacenarán las figuras
\graphicspath{{figures/}}

% =========================================================
% TABLAS
% =========================================================

\usepackage{booktabs}
\usepackage{tabularx}
\usepackage{array}
\usepackage{multirow}
\usepackage{longtable}

% =========================================================
% COLORES
% =========================================================

\usepackage{xcolor}

\definecolor{linkcolor}{RGB}{0,70,140}
\definecolor{citecolor}{RGB}{0,100,70}
\definecolor{urlcolor}{RGB}{120,40,120}

% =========================================================
% CÓDIGO
% =========================================================

\usepackage{listings}

\lstdefinestyle{pythonstyle}{
    language=Python,
    basicstyle=\ttfamily\small,
    numbers=left,
    numberstyle=\tiny,
    stepnumber=1,
    numbersep=8pt,
    frame=single,
    breaklines=true,
    showstringspaces=false,
    tabsize=4,
    captionpos=b
}

\lstset{style=pythonstyle}

% =========================================================
% HIPERVÍNCULOS Y REFERENCIAS CRUZADAS
% =========================================================

\usepackage[
    colorlinks=true,
    linkcolor=linkcolor,
    citecolor=citecolor,
    urlcolor=urlcolor,
    bookmarksopen=true,
    bookmarksnumbered=true
]{hyperref}

\usepackage[nameinlink, noabbrev]{cleveref}

% Nombres utilizados por cleveref
\crefname{equation}{ecuación}{ecuaciones}
\Crefname{equation}{Ecuación}{Ecuaciones}

\crefname{figure}{figura}{figuras}
\Crefname{figure}{Figura}{Figuras}

\crefname{table}{tabla}{tablas}
\Crefname{table}{Tabla}{Tablas}

\crefname{section}{sección}{secciones}
\Crefname{section}{Sección}{Secciones}

\crefname{subsection}{subsección}{subsecciones}
\Crefname{subsection}{Subsección}{Subsecciones}

% =========================================================
% ENCABEZADOS Y PIES DE PÁGINA
% =========================================================

\usepackage{fancyhdr}

\pagestyle{fancy}
\fancyhf{}

\fancyhead[L]{Demostración del efecto del tamaño de dexel}
\fancyhead[R]{\thepage}
\setlength{\headheight}{13.6pt}

\renewcommand{\headrulewidth}{0.4pt}
\renewcommand{\footrulewidth}{0pt}

% =========================================================
% FORMATO DE TÍTULOS
% =========================================================

\usepackage{titlesec}

\titleformat{\section}
    {\Large\bfseries}
    {\thesection.}
    {0.5em}
    {}

\titleformat{\subsection}
    {\large\bfseries}
    {\thesubsection.}
    {0.5em}
    {}

% =========================================================
% ESPACIADO
% =========================================================

\usepackage{setspace}

\setstretch{1.15}
\setlength{\parindent}{0pt}
\setlength{\parskip}{0.6em}

% =========================================================
% BIBLIOGRAFÍA
% =========================================================

\usepackage[
    backend=biber,
    style=numeric,
    sorting=none
]{biblatex}

\addbibresource{references.bib}

% =========================================================
% COMANDOS PERSONALIZADOS
% =========================================================

\newcommand{\ap}{a_p}
\newcommand{\Aref}{A_{\mathrm{ref}}}
\newcommand{\dtlag}{\Delta t}
\newcommand{\dxl}{h_d}
\newcommand{\RMS}{\operatorname{RMS}}
\newcommand{\Real}{\operatorname{Re}}

% =========================================================
% INFORMACIÓN DEL DOCUMENTO
% =========================================================

\title{
    \textbf{Análisis del efecto del tamaño de dexel sobre el crecimiento del chatter}
}

\author{
    Enrique Mireles
}

\date{\today}

% =========================================================
% INICIO DEL DOCUMENTO
% =========================================================

\begin{document}

\maketitle

\begin{abstract}

En este documento se presenta el procedimiento teórico y numérico utilizado
para estudiar el efecto del tamaño de dexel sobre la aparición y el crecimiento
del chatter en simulaciones de mecanizado. El análisis se basa en el ajuste
exponencial de señales dinámicas, la comparación de sus tasas de crecimiento
y la descomposición del retraso temporal observado entre diferentes casos.

\end{abstract}

\tableofcontents

\newpage


\section{Introduction}
\label{sec:introduction}

Dans les simulations étudiées, la géométrie de la pièce est représentée par une
discrétisation en dexels. La taille de dexel, notée \(h_d\), définit la résolution
spatiale du modèle : plus \(h_d\) est faible, plus la discrétisation est fine.

Une différence systématique est observée entre les simulations. Lorsque la taille
de dexel diminue, la partie exponentielle de la réponse dynamique est décalée vers
des temps plus élevés. Autrement dit, le raffinement de la discrétisation semble
retarder le développement exponentiel de la réponse.

Cette observation peut être résumée par

\begin{equation}
    h_{d,j} < h_{d,i}
    \quad\Longrightarrow\quad
    t_{\mathrm{exp},j} > t_{\mathrm{exp},i},
    \label{eq:observation_dexel}
\end{equation}

où \(t_{\mathrm{exp},i}\) désigne l'instant associé à la partie exponentielle du
cas \(i\). Cette relation décrit uniquement l'observation numérique et ne permet
pas encore d'en identifier la cause.

L'objectif de cette étude est de comprendre l'origine de ce décalage temporel.
Deux effets seront examinés :

\begin{enumerate}
    \item une différence du niveau initial de perturbation entre les
    discrétisations ;
    \item une différence du taux de croissance exponentielle.
\end{enumerate}

Le décalage mesuré directement entre deux courbes sera noté

\begin{equation}
    \Delta t_{ij}^{\mathrm{mes}}
    =
    t_j^{\mathrm{mes}}-t_i^{\mathrm{mes}}.
    \label{eq:dt_mes}
\end{equation}

La convention adoptée est donc

\begin{equation}
    \Delta t_{ij}^{\mathrm{mes}}>0
    \quad\Longrightarrow\quad
    \text{le cas \(j\) est retardé par rapport au cas \(i\)}.
\end{equation}

Ce décalage sera ensuite comparé à deux estimations : une première obtenue à
partir des ajustements exponentiels, notée
\(\Delta t_{ij}^{\mathrm{fit}}\), et une seconde obtenue en utilisant la force
comme proxy de la perturbation, notée \(\Delta t_{ij}^{F}\).


% ==============================================
% ======== Seccion 2 ==========================
% ==============================================


\section{Configuration numérique et modèle dynamique}
\label{sec:configuration}

\subsection{Géométrie et profondeur de coupe variable}
\label{subsec:geometrie}

La pièce étudiée est un cône tronqué. Pendant l'usinage, la profondeur de coupe
augmente linéairement de \(a_{p,0}\) à \(a_{p,1}\) :

\begin{equation}
    a_p(t)
    =
    a_{p,0}
    +
    \left(a_{p,1}-a_{p,0}\right)
    \frac{t}{t_{\mathrm{rampe}}},
    \qquad
    0 \leq t \leq t_{\mathrm{rampe}}.
    \label{eq:rampe_ap}
\end{equation}

Dans cette expression, \(a_{p,0}\) et \(a_{p,1}\) sont les profondeurs de coupe
initiale et finale, tandis que \(t_{\mathrm{rampe}}\) est la durée de la rampe.

Dans les simulations considérées, la profondeur de coupe initiale est \(a_{p,0}=\SI{5}{\milli\metre}\) et la profondeur de coupe finale est \(a_{p,1}=\SI{15}{\milli\metre}\).

% Figure à ajouter :
% Schéma du cône tronqué avec la direction d'usinage,
% a_{p,0}, a_{p,1} et la longueur de la zone conique.

\begin{figure}[H]
    \centering
    \fbox{
        \parbox[c][5cm][c]{0.85\textwidth}{
            \centering
            \textbf{Figure à ajouter}\\[0.3cm]
            Schéma du cône tronqué avec la direction d'usinage,
            $a_{p,0}$, $a_{p,1}$ et la longueur de la zone conique.
        }
    }
    \caption{Géométrie de la pièce et évolution de la profondeur de coupe.}
    \label{fig:geometrie_cone}
\end{figure}

% Figure à générer depuis le code :
% Courbe a_p(t) entre 5 mm et 15 mm.

\begin{figure}[H]
    \centering
    \fbox{
        \parbox[c][4.5cm][c]{0.85\textwidth}{
            \centering
            \textbf{Figure à générer depuis le code}\\[0.3cm]
            Courbe $a_p(t)$ entre \SI{5}{\milli\metre}
            et \SI{15}{\milli\metre}.
        }
    }
    \caption{Évolution temporelle de la profondeur de coupe.}
    \label{fig:rampe_ap}
\end{figure}

\subsection{Modèle dynamique}
\label{subsec:modele-dynamique}

La dynamique de l'outil est représentée par un modèle modal à deux degrés de
liberté. Le déplacement de l'outil est projeté selon les directions modales du
système, tandis que l'analyse présentée dans ce document est effectuée sur la
direction axiale.

Pour une direction modale, l'équation du mouvement peut s'écrire sous la forme

\begin{equation}
    m\ddot{x}(t)
    +
    c\dot{x}(t)
    +
    kx(t)
    =
    F_{\mathrm{dyn}}(t),
    \label{eq:modele-dynamique}
\end{equation}

où \(m\), \(c\) et \(k\) représentent respectivement la masse modale,
l'amortissement modal et la rigidité modale. La variable \(x(t)\) désigne le
déplacement dynamique dans la direction étudiée.

% Figure à ajouter :
% Schéma des deux degrés de liberté, directions modales,
% direction de coupe et projection dans la direction axiale.

\begin{figure}[H]
    \centering
    \fbox{
        \parbox[c][5cm][c]{0.85\textwidth}{
            \centering
            \textbf{Figure à ajouter}\\[0.3cm]
            Schéma du modèle à deux degrés de liberté et des directions modales.
        }
    }
    \caption{Modèle dynamique et directions modales considérées.}
    \label{fig:modele_2ddl}
\end{figure}


% \subsection{Modèle régénératif}
% \label{subsec:modele-regeneratif}

L'épaisseur instantanée dépend de la surface laissée au passage précédent :

\begin{equation}
    h(t)=h_0+x(t-T)-x(t),
    \label{eq:epaisseur-regenerative}
\end{equation}

où \(h_0\) est l'épaisseur nominale et \(T\) la période de révolution. La
contribution régénérative est donc donnée par \(x(t-T)-x(t)\).

En supposant une loi de force linéaire, la force dynamique s'écrit

\begin{equation}
    F_{\mathrm{dyn}}(t)
    =
    K_f a_p(t)\left[x(t-T)-x(t)\right],
    \label{eq:force-regenerative}
\end{equation}

où \(K_f\) est le coefficient de force. Le modèle dynamique devient alors

\begin{equation}
    m\ddot{x}(t)+c\dot{x}(t)+kx(t)
    =
    K_f a_p(t)\left[x(t-T)-x(t)\right].
    \label{eq:equation-regenerative}
\end{equation}


\subsection{Position par rapport à la limite de stabilité}
\label{subsec:stabilite}

Pour la vitesse de rotation considérée, le diagramme de stabilité fournit une
profondeur critique \(a_{p,\mathrm{crit}}\).

La rampe étudiée traverse cette limite :

\begin{equation}
    a_{p,0}
    <
    a_{p,\mathrm{crit}}
    <
    a_{p,1}.
    \label{eq:traversee-limite}
\end{equation}

Le système évolue ainsi d'une région située sous la limite théorique vers une
région située au-dessus de celle-ci.

% Figure à ajouter :
% Diagramme des lobes de stabilité avec :
% - la vitesse de rotation étudiée ;
% - la profondeur critique a_{p,crit} ;
% - la trajectoire verticale de 5 mm à 15 mm.

\begin{figure}[H]
    \centering
    \fbox{
        \parbox[c][5cm][c]{0.85\textwidth}{
            \centering
            \textbf{Figure à ajouter}\\[0.3cm]
            Diagramme des lobes de stabilité avec la vitesse de rotation,
            la profondeur critique $a_{p,\mathrm{crit}}$ et la rampe de 5 mm à 15 mm.
        }
    }
    \caption{Position de la rampe de profondeur de coupe dans le diagramme de stabilité.}
    \label{fig:lobes_stabilite}
\end{figure}

\subsection{Signaux analysés}
\label{subsec:signaux}

Pour chaque taille de dexel \(h_{d,i}\), les signaux suivants sont analysés :

\begin{description}
    \item[\(x_i(t)\)] déplacement axial du cas \(i\);
    \item[\(\dot{x}_i(t)\)] vitesse axiale du cas \(i\);
    \item[\(F_i(t)\)] force projetée dans la direction étudiée.
\end{description}

% Figure à générer depuis le code :
% Signaux complets de déplacement, vitesse et force
% pour les différentes tailles de dexel.

\begin{figure}[H]
    \centering
    \fbox{
        \parbox[c][5cm][c]{0.85\textwidth}{
            \centering
            \textbf{Figure à ajouter}\\[0.3cm]
            Signaux complets de déplacement, vitesse et force pour les différentes tailles de dexel.
        }
    }
    \caption{Signaux extraits pour les différentes tailles de dexel.}
    \label{fig:signaux_complets}
\end{figure}

\subsection{Paramètres principaux}
\label{subsec:parametres}

Les principaux paramètres sont regroupés dans la
Tableau~\ref{tab:parametres_simulation}.


\begin{table}[H]
    \centering
    \caption{Principaux paramètres du modèle numérique.}
    \label{tab:parametres_simulation}
    \begin{tabular}{llll}
        \toprule
        Paramètre & Symbole & Valeur & Unité \\
        \midrule

        Vitesse de rotation
        & \(n\)
        & \(12094.28\)
        & \(\si{\per\minute}\) \\

        Avance nominale
        & \(f_z\)
        & \(0.05\)
        & \(\si{\milli\metre}\) \\

        Coefficient de force
        & \(K_f\)
        & \(1.0\times10^{9}\)
        & \(\si{\newton\per\metre\squared}\) \\

        Profondeur initiale
        & \(a_{p,0}\)
        & \(5\)
        & \(\si{\milli\metre}\) \\

        Profondeur finale
        & \(a_{p,1}\)
        & \(15\)
        & \(\si{\milli\metre}\) \\

        Longueur de la zone conique
        & \(L\)
        & \(150\)
        & \(\si{\milli\metre}\) \\

        Rigidité modale
        & \(k\)
        & \(2.13\times10^{8}\)
        & \(\si{\newton\per\metre}\) \\

        Fréquence propre
        & \(f_n\)
        & \(150\)
        & \(\si{\hertz}\) \\

        Pulsation propre
        & \(\omega_n=2\pi f_n\)
        & \(942.48\)
        & \(\si{\radian\per\second}\) \\

        Pulsation propre au carré
        & \(\omega_n^2=k/m\)
        & \(8.8826\times10^{5}\)
        & \(\si{\per\second\squared}\) \\

        Masse modale
        & \(m=k/\omega_n^2\)
        & \(239.79\)
        & \(\si{\kilogram}\) \\

        Angle de projection
        & \(\theta\)
        & \(135\)
        & \(\si{\degree}\) \\

        Coefficient de projection modale
        & \(\phi=\sin(\theta)/\sqrt{m}\)
        & \(4.5663\times10^{-2}\)
        & \(\si{\per\sqrt{\kilogram}}\) \\

        \bottomrule
    \end{tabular}
\end{table}






% ==============================================
% ======== Seccion 3 ==========================
% ==============================================
\section{Prétraitement des signaux}
\label{sec:pretraitement}

La variation de la profondeur de coupe introduit une évolution déterministe dans
la force, le déplacement et la vitesse. Ces contributions sont retirées avant
l'analyse afin de conserver principalement la réponse dynamique.

\subsection{Force nominale et force résiduelle}
\label{subsec:force-residuelle}

La force nominale associée à la profondeur de coupe variable est définie par

\begin{equation}
    F_{\mathrm{nom}}(t)
    =
    K_f\,a_p(t)\,f_z,
    \label{eq:force-nominale}
\end{equation}

avec :

\begin{description}
    \item[\(K_f\)] coefficient de force de coupe ;
    \item[\(a_p(t)\)] profondeur de coupe instantanée ;
    \item[\(f_z\)] avance nominale utilisée dans la simulation.
\end{description}

Cette expression représente l'évolution attendue de la force due uniquement à
la rampe de profondeur de coupe. La force résiduelle du cas \(i\) est alors
obtenue par

\begin{equation}
    \delta F_i(t)
    =
    F_i(t)-F_{\mathrm{nom}}(t),
    \label{eq:force-residuelle}
\end{equation}

où \(F_i(t)\) est la force simulée pour la taille de dexel \(h_{d,i}\).
La quantité \(\delta F_i(t)\) contient les variations qui ne sont pas décrites
par le modèle nominal.

\begin{figure}[H]
    \centering
    \fbox{
        \parbox[c][5cm][c]{0.85\textwidth}{
            \centering
            \textbf{Figure à ajouter}\\[0.3cm]
            Force simulée, force nominale théorique et force résiduelle pour les différentes tailles de dexel.
        }
    }
    \caption{Décomposition de la force simulée en composante nominale et composante résiduelle.}
    \label{fig:force-residuelle}
\end{figure}

\subsection{Correction du déplacement}
\label{subsec:correction-deplacement}

L'augmentation de la force nominale produit une déflexion quasi statique de
l'outil. Dans la formulation modale utilisée, cette déflexion est donnée par

\begin{equation}
    \delta_{\mathrm{stat}}(t)
    =
    \frac{\phi^{2}}{\omega_2}\,
    F_{\mathrm{nom}}(t),
    \label{eq:deflexion-statique}
\end{equation}

avec :

\begin{description}
    \item[\(\phi\)] coefficient de projection modale dans la direction étudiée ;
    \item[\(\omega_2\)] pulsation modale au carré, avec
    \(\omega_2=k/m\).
\end{description}

En substituant la \ref{eq:force-nominale} dans la
\ref{eq:deflexion-statique}, on obtient

\begin{equation}
    \delta_{\mathrm{stat}}(t)
    =
    \frac{\phi^{2}}{\omega_2}\,
    K_f\,a_p(t)\,f_z.
    \label{eq:deflexion-statique-developpee}
\end{equation}

Le déplacement dynamique est ensuite défini par

\begin{equation}
    x_{i,\mathrm{dyn}}(t)
    =
    x_i(t)-\delta_{\mathrm{stat}}(t).
    \label{eq:deplacement-corrige}
\end{equation}

Cette correction retire la tendance imposée par l'augmentation de la force
nominale.

\begin{figure}[H]
    \centering
    \fbox{
        \parbox[c][5cm][c]{0.85\textwidth}{
            \centering
            \textbf{Figure à ajouter}\\[0.3cm]
            Déplacement simulé, déflexion quasi statique théorique et déplacement corrigé pour les différentes tailles de dexel.
        }
    }
    \caption{Correction de la composante quasi statique du déplacement.}
    \label{fig:deplacement-corrige}
\end{figure}

\subsection{Correction de la vitesse}
\label{subsec:correction-vitesse}

La composante déterministe de la vitesse est obtenue en dérivant la déflexion
statique :

\begin{equation}
    \dot{\delta}_{\mathrm{stat}}(t)
    =
    \frac{d\delta_{\mathrm{stat}}(t)}{dt}.
    \label{eq:vitesse-statique}
\end{equation}

À partir de la \ref{eq:deflexion-statique-developpee}, il vient

\begin{equation}
    \dot{\delta}_{\mathrm{stat}}(t)
    =
    \frac{\phi^{2}}{\omega_2}\,
    K_f\,f_z\,
    \frac{da_p(t)}{dt}.
    \label{eq:vitesse-statique-developpee}
\end{equation}

Comme \(a_p(t)\) est linéaire pendant la rampe,

\begin{equation}
    \frac{da_p(t)}{dt}
    =
    \frac{a_{p,1}-a_{p,0}}{t_{\mathrm{rampe}}},
    \qquad
    0\leq t\leq t_{\mathrm{rampe}}.
    \label{eq:derivee-ap}
\end{equation}

La vitesse dynamique corrigée est alors

\begin{equation}
    \dot{x}_{i,\mathrm{dyn}}(t)
    =
    \dot{x}_i(t)-\dot{\delta}_{\mathrm{stat}}(t).
    \label{eq:vitesse-corrigee}
\end{equation}

\begin{figure}[H]
    \centering
    \fbox{
        \parbox[c][5cm][c]{0.85\textwidth}{
            \centering
            \textbf{Figure à ajouter}\\[0.3cm]
            Vitesse simulée, vitesse associée à la déflexion statique et vitesse corrigée pour les différentes tailles de dexel.
        }
    }
    \caption{Correction de la composante déterministe de la vitesse.}
    \label{fig:vitesse-corrigee}
\end{figure}

\subsection{Normalisation de la force résiduelle}
\label{subsec:normalisation-force}

D'après le modèle régénératif, la force dynamique contient directement le
facteur \(a_p(t)\) :

\begin{equation}
    \delta F_i(t)
    \approx
    K_f\,a_p(t)
    \left[
        x_i(t-T)-x_i(t)
    \right].
    \label{eq:force-residuelle-regenerative}
\end{equation}

Deux perturbations évaluées à des instants différents peuvent donc présenter
des amplitudes différentes uniquement parce que la profondeur de coupe a
évolué. Pour supprimer cet effet direct, la force résiduelle est normalisée par
\(a_p(t)\) :

\begin{equation}
    \delta F_{i,\mathrm{norm}}(t)
    =
    \frac{\delta F_i(t)}{a_p(t)}.
    \label{eq:force-normalisee}
\end{equation}

En utilisant la \ref{eq:force-residuelle-regenerative}, on obtient

\begin{equation}
    \delta F_{i,\mathrm{norm}}(t)
    \approx
    K_f
    \left[
        x_i(t-T)-x_i(t)
    \right].
    \label{eq:force-normalisee-regenerative}
\end{equation}

La normalisation retire ainsi la dépendance directe à \(a_p(t)\). Cette quantité
sera utilisée plus loin comme proxy de la perturbation associée à la
discrétisation.

\begin{figure}[H]
    \centering
    \fbox{
        \parbox[c][5cm][c]{0.85\textwidth}{
            \centering
            \textbf{Figure à ajouter}\\[0.3cm]
            Comparaison de la force résiduelle et de la force résiduelle normalisée pour les différentes tailles de dexel.
        }
    }
    \caption{Effet de la normalisation de la force résiduelle par la profondeur de coupe.}
    \label{fig:force-normalisee}
\end{figure}


%==============================================
% ======= ===== Seccion 4 ==========================
%==============================================

\section{Construction de la grandeur d'amplitude}
\label{sec:amplitude}

L'analyse porte sur l'évolution de l'amplitude des signaux corrigés. Afin de
réduire les oscillations rapides tout en conservant leur évolution globale, une
valeur efficace glissante est calculée.

\subsection{RMS glissant}
\label{subsec:rms-glissant}

Pour un signal corrigé \(z_i(t)\), le RMS glissant est défini par

\begin{equation}
    R_i(t)
    =
    \sqrt{
        \frac{1}{T_w}
        \int_{t-T_w/2}^{t+T_w/2}
        z_i^2(\tau)\,d\tau
    },
    \label{eq:rms-continu}
\end{equation}

avec :

\begin{description}
    \item[\(R_i(t)\)] amplitude RMS du cas \(i\) ;
    \item[\(z_i(t)\)] signal corrigé analysé ;
    \item[\(T_w\)] durée de la fenêtre glissante.
\end{description}

Dans la mise en œuvre numérique, cette expression est approchée par

\begin{equation}
    R_i(t_k)
    =
    \sqrt{
        \frac{1}{N_w}
        \sum_{n\in\mathcal{W}_k}
        z_i^2(t_n)
    },
    \label{eq:rms-discret}
\end{equation}

où \(N_w\) est le nombre d'échantillons contenus dans la fenêtre
\(\mathcal{W}_k\).

La fenêtre utilisée couvre plusieurs périodes de la réponse dominante :

\begin{equation}
    T_w
    =
    \frac{5}{f_n},
    \label{eq:fenetre-rms}
\end{equation}

où \(f_n=\SI{150}{\hertz}\) est la fréquence propre considérée. Ce choix permet
d'obtenir une enveloppe suffisamment lisse sans supprimer l'évolution
temporelle du signal.

\begin{figure}[H]
    \centering
    \fbox{
        \parbox[c][5cm][c]{0.85\textwidth}{
            \centering
            \textbf{Figure à ajouter}\\[0.3cm]
            Signal corrigé et RMS glissant pour une taille de dexel représentative.
        }
    }
    \caption{Construction de l'amplitude RMS à partir du signal corrigé.}
    \label{fig:rms-construction}
\end{figure}

\subsection{Représentation logarithmique}
\label{subsec:representation-log}

Dans un intervalle où la réponse suit une évolution exponentielle,

\begin{equation}
    R_i(t)
    =
    P_i\,\exp\left(\beta_i t\right),
    \label{eq:rms-exponentiel}
\end{equation}

le logarithme naturel transforme cette relation en une droite :

\begin{align}
    \ln R_i(t)
    &=
    \ln\left[
        P_i\,\exp\left(\beta_i t\right)
    \right],\\
    \ln R_i(t)
    &=
    \ln P_i+\beta_i t.
    \label{eq:passage-log}
\end{align}

En posant

\begin{equation}
    c_i=\ln P_i,
    \label{eq:def-c}
\end{equation}

le modèle ajusté devient

\begin{equation}
    \ln R_i(t)
    =
    c_i+\beta_i t.
    \label{eq:modele-log}
\end{equation}

Toutes les courbes utilisées dans la suite sont donc étudiées en échelle
logarithmique naturelle.

\begin{description}
    \item[\(c_i\)] ordonnée à l'origine de la droite ajustée ;
    \item[\(P_i=\exp(c_i)\)] amplitude extrapolée au temps \(t=0\) ;
    \item[\(\beta_i\)] taux de croissance exponentielle, en
    \(\si{\per\second}\).
\end{description}

\subsection{Sélection des intervalles d'étude}
\label{subsec:intervalles-etude}

Chaque taille de dexel possède son propre intervalle d'étude,
\(I_i=[t_{a,i},t_{b,i}]\). Cet intervalle est choisi dans la zone où
\(\ln R_i(t)\) présente un comportement approximativement linéaire.

L'ajustement est réalisé uniquement dans cet intervalle :

\begin{equation}
    \ln R_i(t)
    \approx
    c_i+\beta_i t,
    \qquad
    t\in I_i.
    \label{eq:fit-intervalle}
\end{equation}

Ainsi, le modèle exponentiel n'est pas supposé valable sur toute la simulation,
mais seulement dans la région étudiée.

\begin{figure}[H]
    \centering
    \fbox{
        \parbox[c][5cm][c]{0.85\textwidth}{
            \centering
            \textbf{Figure à ajouter}\\[0.3cm]
            Courbes \(\ln R_i(t)\) pour les différentes tailles de dexel,
            avec les intervalles d'ajustement propres à chaque cas.
        }
    }
    \caption{Sélection des intervalles présentant un comportement log-linéaire.}
    \label{fig:intervalles-log}
\end{figure}

\subsection{Qualité de l'ajustement}
\label{subsec:qualite-ajustement}

La qualité de l'ajustement linéaire est évaluée avec le coefficient de
détermination \(R_i^2\). Pour les valeurs observées
\(y_k=\ln R_i(t_k)\) et les valeurs ajustées
\(\widehat{y}_k=c_i+\beta_i t_k\), il est défini par

\begin{equation}
    R_i^2
    =
    1-
    \frac{
        \displaystyle\sum_k
        \left(y_k-\widehat{y}_k\right)^2
    }{
        \displaystyle\sum_k
        \left(y_k-\overline{y}\right)^2
    },
    \label{eq:r2}
\end{equation}

où \(\overline{y}\) est la moyenne des valeurs observées.

Un coefficient \(R_i^2\) proche de l'un indique que l'évolution de
\(\ln R_i(t)\) est correctement représentée par une droite dans l'intervalle
sélectionné. Il constitue un critère de qualité de l'ajustement, sans démontrer
à lui seul la validité physique du modèle exponentiel.

\begin{figure}[H]
    \centering
    \fbox{
        \parbox[c][5cm][c]{0.85\textwidth}{
            \centering
            \textbf{Figure à ajouter}\\[0.3cm]
            Courbes \(\ln R_i(t)\), droites ajustées, valeurs de \(\beta_i\)
            et coefficients \(R_i^2\).
        }
    }
    \caption{Ajustements linéaires des amplitudes en échelle logarithmique.}
    \label{fig:fits-log}
\end{figure}



% =============================================
% === seccion 5 ==========================
% =============================================
\section{Niveau de référence et définition des décalages temporels}
\label{sec:niveau-reference}

L'objectif est de comparer les différentes courbes à un même niveau
d'amplitude. Un niveau de référence \(A_{\mathrm{ref}}\) est donc choisi de
manière à être atteint par toutes les courbes dans leurs intervalles d'étude
respectifs.

En échelle logarithmique, cette condition s'écrit

\begin{equation}
    \ln A_{\mathrm{ref}}
    \in
    \bigcap_{i=1}^{N}
    \left[
        \min_{t\in I_i}\ln R_i(t),
        \max_{t\in I_i}\ln R_i(t)
    \right].
    \label{eq:condition-Aref}
\end{equation}

Le niveau \(A_{\mathrm{ref}}\) n'est pas un seuil physique. Il sert uniquement
à définir une référence commune pour mesurer les décalages temporels.

\begin{figure}[H]
    \centering
    \fbox{
        \parbox[c][5cm][c]{0.85\textwidth}{
            \centering
            \textbf{Figure à ajouter}\\[0.3cm]
            Courbes \(\ln R_i(t)\) pour les différentes tailles de dexel,
            avec la droite horizontale \(\ln A_{\mathrm{ref}}\) et les temps de croisement.
        }
    }
    \caption{Définition du niveau de référence commun.}
    \label{fig:niveau-reference}
\end{figure}

\subsection{Décalage mesuré directement}
\label{subsec:decalage-mesure}

Pour chaque cas \(i\), le temps \(t_i^{\mathrm{mes}}\) est défini comme
l'instant où la courbe mesurée atteint \(A_{\mathrm{ref}}\) :

\begin{equation}
    R_i\!\left(t_i^{\mathrm{mes}}\right)
    =
    A_{\mathrm{ref}}.
    \label{eq:temps-mesure}
\end{equation}

Pour deux cas \(i\) et \(j\), le décalage mesuré est

\begin{equation}
    \Delta t_{ij}^{\mathrm{mes}}
    =
    t_j^{\mathrm{mes}}
    -
    t_i^{\mathrm{mes}}.
    \label{eq:decalage-mesure}
\end{equation}

La convention de signe adoptée est la suivante :

\begin{description}
    \item[\(\Delta t_{ij}^{\mathrm{mes}}>0\)] le cas \(j\) atteint
    \(A_{\mathrm{ref}}\) après le cas \(i\) ;
    \item[\(\Delta t_{ij}^{\mathrm{mes}}<0\)] le cas \(j\) atteint
    \(A_{\mathrm{ref}}\) avant le cas \(i\).
\end{description}

Dans l'implémentation numérique, \(t_i^{\mathrm{mes}}\) est obtenu à partir du
premier croisement de la courbe RMS avec \(A_{\mathrm{ref}}\).


\subsection{Décalage obtenu par l’ajustement}
\label{subsec:decalage-fit}

Dans l’intervalle d’étude \(I_i\), la courbe logarithmique est décrite par

\begin{equation}
    \ln R_i(t)
    =
    c_i+\beta_i t
    =
    \ln P_i+\beta_i t,
    \label{eq:modele-log-fit}
\end{equation}

où \(P_i=\exp(c_i)\) est l’amplitude initiale extrapolée à \(t=0\) à partir de
l’ajustement.

Le temps auquel la droite ajustée atteint le niveau \(A_{\mathrm{ref}}\) vérifie

\begin{equation}
    \ln A_{\mathrm{ref}}
    =
    c_i^{\mathrm{fit}}+\beta_i^{\mathrm{fit}} t_i^{\mathrm{fit}}
    =
    \ln P_i^{\mathrm{fit}}+\beta_i^{\mathrm{fit}} t_i^{\mathrm{fit}}.
    \label{eq:condition-fit-Aref}
\end{equation}

En isolant \(t_i^{\mathrm{fit}}\), on obtient

\begin{align}
    t_i^{\mathrm{fit}}
    &=
    \frac{\ln A_{\mathrm{ref}}-c_i^{\mathrm{fit}}}{\beta_i^{\mathrm{fit}}},
    \nonumber\\
    &=
    \frac{\ln A_{\mathrm{ref}}-\ln P_i^{\mathrm{fit}}}{\beta_i^{\mathrm{fit}}},
    \nonumber\\
    &=
    \ln\left(
        \frac{1}{\beta_i^{\mathrm{fit}}}
        \frac{A_{\mathrm{ref}}}{P_i^{\mathrm{fit}}}
    \right).
    \label{eq:temps-fit}
\end{align}

Pour deux cas \(i\) et \(j\), le décalage prédit par les ajustements est défini
par

\begin{equation}
    \Delta t_{ij}^{\mathrm{fit}}
    =
    t_j^{\mathrm{fit}}
    -
    t_i^{\mathrm{fit}}.
    \label{eq:decalage-fit}
\end{equation}

En substituant la \ref{eq:temps-fit}, il vient

\begin{align}
    \Delta t_{ij}^{\mathrm{fit}}
    &=
    \frac{\ln A_{\mathrm{ref}}-c_j^{\mathrm{fit}}}{\beta_j^{\mathrm{fit}}}
    -
    \frac{\ln A_{\mathrm{ref}}-c_i^{\mathrm{fit}}}{\beta_i^{\mathrm{fit}}},
    \nonumber\\
    &=
    \frac{\ln A_{\mathrm{ref}}-\ln P_j^{\mathrm{fit}}}{\beta_j^{\mathrm{fit}}}
    -
    \frac{\ln A_{\mathrm{ref}}-\ln P_i^{\mathrm{fit}}}{\beta_i^{\mathrm{fit}}},
    \nonumber\\
    &=
    \frac{1}{\beta_j^{\mathrm{fit}}}
    \ln\left(
        \frac{A_{\mathrm{ref}}}{P_j^{\mathrm{fit}}}
    \right)
    -
    \frac{1}{\beta_i^{\mathrm{fit}}}
    \ln\left(
        \frac{A_{\mathrm{ref}}}{P_i^{\mathrm{fit}}}
    \right).
    \label{eq:decalage-fit-general}
\end{align}


\subsection{Décalage estimé à partir de la force}
\label{subsec:decalage-force}

L’amplitude \(P_i^{\mathrm{fit}}=\exp(c_i^{\mathrm{fit}})\) obtenue par l’ajustement est une valeur extrapolée
à \(t=0\). Afin d’estimer cette amplitude indépendamment du fitting, la force
résiduelle normalisée est utilisée comme proxy de la perturbation associée à la
discrétisation.


\subsubsection{Définition du proxy de perturbation}
\label{subsubsec:relation-proxy-amplitude}

La force résiduelle normalisée du cas \(i\) est donnée par l'équation \ref{eq:force-normalisee-regenerative}. Le proxy de perturbation est défini comme la valeur RMS locale de cette force
à l’instant \(t_{s,i}\) :

\begin{equation}
    P_{h,i}
    =
    \operatorname{RMS}
    \left[
        \delta F_{i,\mathrm{norm}}(t)
    \right]_{t=t_{s,i}}.
    \label{eq:definition-proxy-force}
\end{equation}

En substituant la \ref{eq:force-normalisee-regenerative}, on obtient

\begin{equation}
    P_{h,i}
    =
    K_f
    \operatorname{RMS}
    \left[
        x_i(t-T)-x_i(t)
    \right]_{t=t_{s,i}}.
    \label{eq:proxy-difference-regenerative}
\end{equation}

Dans l’intervalle d’ajustement, le déplacement est supposé dominé par une
réponse oscillatoire dont l’enveloppe évolue exponentiellement. Le déplacement
\(x_i(t)\) et sa version retardée \(x_i(t-T)\) possèdent donc le même taux de
croissance \(\beta_i\). Leur différence conserve également cette enveloppe
exponentielle.

La valeur RMS de la différence régénérative peut alors être écrite sous la
forme

\begin{equation}
    \operatorname{RMS}
    \left[
        x_i(t-T)-x_i(t)
    \right]
    \approx
    C_i R_i^{\mathrm{F}}(t),
    \label{eq:proportion-difference-deplacement}
\end{equation}

où \(C_i\) est un facteur de proportionnalité qui dépend principalement du
retard \(T\), de la fréquence dominante et de la relation de phase entre
\(x_i(t-T)\) et \(x_i(t)\).

Le proxy de force devient donc

\begin{equation}
    P_{h,i}
    \approx
    K_f C_i R_i^{\mathrm{F}}(t_{s,i}).
    \label{eq:proxy-amplitude-simplifie}
\end{equation}

Cette relation ne signifie pas que la force normalisée est égale au
déplacement. Elle indique uniquement que, dans l’intervalle exponentiel, les
deux grandeurs possèdent la même évolution d’amplitude, à un facteur
multiplicatif près.

Pour comparer directement les amplitudes initiales entre plusieurs tailles de
dexel, le facteur \(C_i\) doit rester identique ou suffisamment proche entre les
cas étudiés. Cette hypothèse pourra être vérifiée à partir de la fréquence
dominante des signaux.

\subsubsection{Expression analytique de \(C_i\)}
\label{subsubsec:expression-analytique-Ci}

Le facteur \(C_i\) peut être calculé explicitement à partir des paramètres du
système. Dans la zone de croissance exponentielle, le déplacement dynamique est
de la forme

\begin{equation}
    x_i(t) = A_i\,e^{\beta_i t}\sin(\omega_i t + \varphi_i),
    \label{eq:signal-exponentiel-sinusoidal}
\end{equation}

où \(\omega_i\) est la pulsation dominante du chatter et \(\varphi_i\) une phase
initiale. Le signal retardé s'écrit

\begin{equation}
    x_i(t-T) = A_i\,e^{\beta_i(t-T)}\sin(\omega_i t - \omega_i T + \varphi_i).
\end{equation}

La différence régénérative est alors

\begin{equation}
    x_i(t-T) - x_i(t)
    =
    A_i\,e^{\beta_i t}
    \bigl[
        e^{-\beta_i T}\sin(\omega_i t - \omega_i T + \varphi_i)
        -
        \sin(\omega_i t + \varphi_i)
    \bigr].
\end{equation}

En calculant la valeur RMS de cette expression sur plusieurs cycles (en
supposant \(\beta_i T \ll 1\), ce qui est généralement satisfait près de la
limite de stabilité), on obtient

\begin{equation}
    \operatorname{RMS}\left[x_i(t-T)-x_i(t)\right]
    \approx
    \frac{A_i\,e^{\beta_i t}}{\sqrt{2}}
    \sqrt{1 + e^{-2\beta_i T} - 2\,e^{-\beta_i T}\cos(\omega_i T)},
\end{equation}

et comme \(\operatorname{RMS}[x_i(t)] = A_i e^{\beta_i t}/\sqrt{2}\), le
facteur de proportionnalité s'exprime sous la forme

\begin{equation}
    \boxed{
        C_i
        =
        \sqrt{1 + e^{-2\beta_i T} - 2\,e^{-\beta_i T}\cos(\omega_i T)}.
    }
    \label{eq:Ci-analytique}
\end{equation}

Ce facteur dépend uniquement du retard régénératif \(T\), du taux de croissance
\(\beta_i\) et de la pulsation dominante \(\omega_i\).

\paragraph{Évaluation numérique.}
Dans les simulations considérées, la vitesse de rotation est telle que
\(T = 60/(n_{\mathrm{rpm}}\cdot N_z)\) avec \(N_z=1\) dent. La pulsation de
chatter est proche de la pulsation modale \(\omega_i \approx 2\pi\times
\SI{150}{\hertz}\). Puisque \(\beta_i T \ll 1\) (le taux de croissance est de
l'ordre de quelques \(\si{\per\second}\) et \(T\) de l'ordre de la
milliseconde), le terme \(e^{-\beta_i T}\approx 1\) et l'expression se
simplifie en

\begin{equation}
    C_i
    \approx
    \sqrt{2\bigl(1 - \cos(\omega_i T)\bigr)}
    =
    2\left|\sin\!\left(\frac{\omega_i T}{2}\right)\right|.
    \label{eq:Ci-approx}
\end{equation}

Comme \(\omega_i\) est très proche d'une pulsation propre du système et que
cette pulsation est commune à tous les cas étudiés (même modèle dynamique,
même vitesse de rotation), on a \(\omega_i \approx \omega_j\) et donc

\begin{equation}
    C_i \approx C_j,
    \label{eq:Ci-Cj-approx}
\end{equation}

ce qui justifie l'hypothèse utilisée dans les sections suivantes.


\subsubsection{Correction temporelle du proxy}
\label{subsubsec:correction-temporelle-proxy}

Le proxy \(P_{h,i}\) est évalué à l’instant \(t_{s,i}\), alors que l’amplitude
$P_i^{\mathrm{fit}}$ du modèle exponentiel correspond à une valeur extrapolée à \(t=0\).
Ces deux grandeurs ne peuvent donc pas être comparées directement.

Dans l’intervalle d’ajustement, l’amplitude dynamique est modélisée par

\begin{equation}
    R_i(t)
    =
    P_i\exp\left(\beta_i t\right),
    \label{eq:modele-exp-correction-proxy}
\end{equation}


À l’instant \(t_{s,i}\),

\begin{equation}
    R_i^{\mathrm{F}}(t_{s,i})
    =
    P_i^{\mathrm{F}}\exp\left(\beta_i^{\mathrm{F}} t_{s,i}\right).
    \label{eq:amplitude-ts-correction}
\end{equation}

En substituant cette expression dans la
\eqref{eq:proxy-amplitude-simplifie}, on obtient

\begin{equation}
    P_{h,i}
    \approx
    K_f C_i P_i^{\mathrm{F}}
    \exp\left(\beta_i^{\mathrm{F}} t_{s,i}\right).
    \label{eq:proxy-avec-croissance}
\end{equation}

Le terme

\begin{equation}
    \exp\left(\beta_i^{\mathrm{F}} t_{s,i}\right)
\end{equation}

représente la croissance accumulée entre l’origine temporelle \(t=0\) et
l’instant d’évaluation \(t_{s,i}\). Il doit donc être supprimé afin de
reconstruire une amplitude équivalente à l’origine temporelle.

En isolant \(P_i^{\mathrm{F}}\), il vient

\begin{align}
    P_i^{\mathrm{F}}
    &=
    \frac{P_{h,i}}
    {K_f C_i
    \exp\left(\beta_i^{\mathrm{F}} t_{s,i}\right)},
    \nonumber\\
    &=
    \frac{P_{h,i}}{K_f C_i}
    \exp\left(-\beta_i^{\mathrm{F}} t_{s,i}\right).
    \label{eq:Pi-reconstruit-force}
\end{align}

L'expression \eqref{eq:Pi-reconstruit-force} est l’amplitude initiale reconstruite à partir de la force.


Cette correction est nécessaire car deux valeurs \(P_{h,i}\) et \(P_{h,j}\)
peuvent être différentes simplement parce qu’elles ont été évaluées à des
instants \(t_{s,i}\) et \(t_{s,j}\) différents. Sans correction temporelle,
cette différence pourrait être attribuée à tort à la taille de dexel.

En échelle logarithmique, la relation
\ref{eq:proxy-avec-croissance} devient

\begin{align}
    \ln P_{h,i}
    &\approx
    \ln K_f
    +
    \ln C_i
    +
    \ln P_i^{\mathrm{F}}
    +
    \beta_i^{\mathrm{F}} t_{s,i},
    \nonumber\\
    &\approx
    \ln K_f
    +
    \ln C_i
    +
    c_i^{\mathrm{F}}
    +
    \beta_i^{\mathrm{F}} t_{s,i},
    \label{eq:log-proxy-correction-temporelle}
\end{align}

avec

\begin{equation}
    c_i^{\mathrm{F}}=\ln P_i^{\mathrm{F}}.
\end{equation}

L’ordonnée à l’origine reconstruite à partir du proxy de force est alors

\begin{equation}
    c_i^{\mathrm{F}}
    =
    \ln P_{h,i}
    -
    \ln K_f
    -
    \ln C_i
    -
    \beta_i^{\mathrm{F}} t_{s,i}.
    \label{eq:ci-force-corrige}
\end{equation}

Si le facteur \(C_i\) peut être considéré comme identique pour tous les cas,
il se simplifiera ensuite dans le calcul des décalages relatifs.




\subsubsection{Temps et décalage estimés}

Le temps auquel le modèle reconstruit à partir de la force atteint
\(A_{\mathrm{ref}}\) est défini par

\begin{equation}
    \ln A_{\mathrm{ref}}
    =
    c_i^{F}
    +
    \beta_i^{\mathrm{F}} t_i^{\mathrm{F}}.
    \label{eq:condition-force-Aref}
\end{equation}

En isolant \(t_i^{F}\),

\begin{align}
    t_i^{F}
    &=
    \frac{\ln A_{\mathrm{ref}}-c_i^{\mathrm{F}}}{\beta_i^{\mathrm{F}}},
    \nonumber\\
    &=
    \frac{
        \ln A_{\mathrm{ref}}
        -
        \ln P_{h,i}
        +
        \beta_i^{\mathrm{F}} t_{s,i}
        +
        \ln K_f
    }{\beta_i},
    \nonumber\\
    &=
    t_{s,i}
    +
    \frac{1}{\beta_i^{\mathrm{F}}}
    \ln\left(
        \frac{K_f A_{\mathrm{ref}}}{P_{h,i}}
    \right).
    \label{eq:temps-force}
\end{align}

Pour deux cas \(i\) et \(j\), le décalage estimé à partir de la force est

\begin{equation}
    \Delta t_{ij}^{F}
    =
    t_j^{F}-t_i^{F}.
    \label{eq:decalage-force}
\end{equation}

En substituant la \ref{eq:temps-force}, on obtient l’expression générale

\begin{align}
    \Delta t_{ij}^{F}
    &=
    t_{s,j}-t_{s,i}
    +
    \frac{1}{\beta_j^{\mathrm{F}}}
    \ln\left(
        \frac{K_f A_{\mathrm{ref}}}{P_{h,j}}
    \right)
    -
    \frac{1}{\beta_i^{\mathrm{F}}}
    \ln\left(
        \frac{K_f A_{\mathrm{ref}}}{P_{h,i}}
    \right).
    \label{eq:decalage-force-general}
\end{align}

Cette expression sera ensuite simplifiée pour l’hypothèse
\(\beta_i=\beta_j\), puis décomposée dans le cas général
\(\beta_i\neq\beta_j\).

\begin{figure}[H]
    \centering
    \fbox{
        \parbox[c][5cm][c]{0.85\textwidth}{
            \centering
            \textbf{Figure à ajouter}\\[0.3cm]
            RMS de la force résiduelle normalisée pour les différentes tailles
            de dexel, avec les instants \(t_{s,i}\) et les valeurs \(P_{h,i}\).
        }
    }
    \caption{Évaluation du proxy de perturbation à partir de la force normalisée.}
    \label{fig:proxy-force}
\end{figure}

\subsection{Comparaison des décalages}
\label{subsec:comparaison-decalages}

L’analyse repose désormais sur trois décalages complémentaires :

\begin{description}
    \item[\(\Delta t_{ij}^{\mathrm{mes}}\)] décalage obtenu directement à partir
    du croisement des courbes avec \(A_{\mathrm{ref}}\) ;

    \item[\(\Delta t_{ij}^{\mathrm{fit}}\)] décalage calculé à partir des
    paramètres \(P_i\) et \(\beta_i\) des ajustements ;

    \item[\(\Delta t_{ij}^{F}\)] décalage reconstruit à partir du proxy
    \(P_{h,i}\), après correction temporelle vers \(t=0\).
\end{description}

Le décalage mesuré constitue la référence. La comparaison avec
\(\Delta t_{ij}^{\mathrm{fit}}\) permet d’évaluer la représentation
exponentielle, tandis que la comparaison avec \(\Delta t_{ij}^{F}\) permet de
tester si la force normalisée explique les différences d’amplitude initiale
associées à la discrétisation.

\begin{figure}[H]
    \centering
    \fbox{
        \parbox[c][5cm][c]{0.85\textwidth}{
            \centering
            \textbf{Figure à ajouter}\\[0.3cm]
            Comparaison de \(\Delta t_{ij}^{\mathrm{mes}}\),
            \(\Delta t_{ij}^{\mathrm{fit}}\) et
            \(\Delta t_{ij}^{F}\) pour chaque paire de tailles de dexel.
        }
    }
    \caption{Comparaison des décalages mesuré, ajusté et reconstruit à partir de la force.}
    \label{fig:comparaison-trois-decalages}
\end{figure}


% =============================================
% === seccion 6 ==========================
% =============================================

\section{Première hypothèse : taux de croissance identiques}
\label{sec:beta-identiques}

La première hypothèse consiste à considérer que la taille de dexel modifie le
niveau initial de la perturbation, mais pas le taux de croissance identifié dans
la zone exponentielle.

Pour deux cas \(i\) et \(j\), cette hypothèse s’écrit

\begin{equation}
    \beta_i^{\mathrm{fit}}
    =
    \beta_j^{\mathrm{fit}}
    =
    \beta^{\mathrm{fit}}.
    \label{eq:hypothese-beta-identiques}
\end{equation}

En échelle logarithmique, les droites ajustées possèdent alors la même pente.
Le décalage entre les deux cas dépend donc uniquement de leurs niveaux initiaux
extrapolés.

\subsection{Simplification du décalage obtenu par ajustement}
\label{subsec:simplification-dt-fit-beta-identiques}

Le décalage général obtenu par ajustement a été défini à la
\ref{eq:decalage-fit-general}. En imposant la
\ref{eq:hypothese-beta-identiques}, il devient

\begin{align}
    \Delta t_{ij}^{\mathrm{fit}}
    &=
    \frac{\ln A_{\mathrm{ref}}-\ln P_j^{\mathrm{fit}}}
    {\beta^{\mathrm{fit}}}
    -
    \frac{\ln A_{\mathrm{ref}}-\ln P_i^{\mathrm{fit}}}
    {\beta^{\mathrm{fit}}},
    \nonumber\\
    &=
    \frac{1}{\beta^{\mathrm{fit}}}
    \left[
        \ln P_i^{\mathrm{fit}}
        -
        \ln P_j^{\mathrm{fit}}
    \right],
    \nonumber\\
    &=
    \frac{1}{\beta^{\mathrm{fit}}}
    \ln\left(
        \frac{P_i^{\mathrm{fit}}}
        {P_j^{\mathrm{fit}}}
    \right).
    \label{eq:dt-fit-beta-identiques}
\end{align}

Comme

\begin{equation}
    c_i^{\mathrm{fit}}=\ln P_i^{\mathrm{fit}},
    \qquad
    c_j^{\mathrm{fit}}=\ln P_j^{\mathrm{fit}},
\end{equation}

la même expression peut s’écrire

\begin{equation}
    \Delta t_{ij}^{\mathrm{fit}}
    =
    \frac{
        c_i^{\mathrm{fit}}-c_j^{\mathrm{fit}}
    }{
        \beta^{\mathrm{fit}}
    }
    =
    \frac{
        \ln P_i^{\mathrm{fit}}
        -
        \ln P_j^{\mathrm{fit}}
    }{
        \beta^{\mathrm{fit}}
    }.
    \label{eq:dt-fit-beta-identiques-ci}
\end{equation}

Le niveau \(A_{\mathrm{ref}}\) disparaît de l’expression. Lorsque les deux
ajustements ont exactement la même pente, leur séparation horizontale est donc
constante, quel que soit le niveau de référence utilisé.

\subsection{Simplification du décalage reconstruit à partir de la force}
\label{subsec:simplification-dt-force-beta-identiques}

Le décalage général reconstruit à partir de la force a été établi à la
\ref{eq:decalage-force-general}. Sous l’hypothèse
\(\beta_i^{\mathrm{fit}}=\beta_j^{\mathrm{fit}}
=\beta^{\mathrm{fit}}\), cette expression devient

\begin{align}
    \Delta t_{ij}^{F}
    &=
    t_{s,j}-t_{s,i} +
    \frac{1}{\beta^{\mathrm{fit}}}
    \ln\left(
        \frac{K_f C_j A_{\mathrm{ref}}}
        {P_{h,j}^{F}}
    \right)
    -
    \frac{1}{\beta^{\mathrm{fit}}}
    \ln\left(
        \frac{K_f C_i A_{\mathrm{ref}}}
        {P_{h,i}^{F}}
    \right).
    \label{eq:dt-force-beta-identiques-intermediaire}
\end{align}

En regroupant les logarithmes,

\begin{align}
    \Delta t_{ij}^{F}
    &=
    t_{s,j}-t_{s,i}
    +
    \frac{1}{\beta^{\mathrm{fit}}}
    \ln\left[
        \frac{
            K_f C_j A_{\mathrm{ref}}
        }{
            P_{h,j}^{F}
        }
        \frac{
            P_{h,i}^{F}
        }{
            K_f C_i A_{\mathrm{ref}}
        }
    \right],
    \nonumber\\
    &=
    t_{s,j}-t_{s,i}
    +
    \frac{1}{\beta^{\mathrm{fit}}}
    \ln\left(
        \frac{P_{h,i}^{F}}{P_{h,j}^{F}}
        \frac{C_j}{C_i}
    \right).
    \label{eq:dt-force-beta-identiques-general}
\end{align}

Le coefficient de force \(K_f\) et le niveau \(A_{\mathrm{ref}}\) disparaissent,
car ils sont communs aux deux cas.

Si le facteur reliant la différence régénérative au déplacement peut être
considéré comme identique,

\begin{equation}
    C_i \approx C_j,
    \label{eq:hypothese-C-identiques}
\end{equation}

alors

\begin{equation}
    \Delta t_{ij}^{F}
    \approx
    t_{s,j}-t_{s,i}
    +
    \frac{1}{\beta^{\mathrm{fit}}}
    \ln\left(
        \frac{P_{h,i}^{F}}
        {P_{h,j}^{F}}
    \right).
    \label{eq:dt-force-beta-identiques-simplifie}
\end{equation}

Le premier terme corrige la différence entre les instants d’évaluation des
proxies. Le second terme représente le décalage associé au rapport entre les
niveaux de perturbation mesurés à partir de la force.

\subsection{Évaluation de l’hypothèse}
\label{subsec:evaluation-beta-identiques}

L’hypothèse de taux identiques est évaluée à partir des trois décalages déjà
définis à la \ref{sec:niveau-reference} :

\begin{description}
    \item[\(\Delta t_{ij}^{\mathrm{mes}}\)]
    décalage mesuré directement sur les courbes ;

    \item[\(\Delta t_{ij}^{\mathrm{fit}}\)]
    décalage prédit par les amplitudes initiales issues des ajustements ;

    \item[\(\Delta t_{ij}^{F}\)]
    décalage reconstruit à partir du proxy de force.
\end{description}

Une proximité entre
\(\Delta t_{ij}^{\mathrm{mes}}\) et
\(\Delta t_{ij}^{\mathrm{fit}}\) indiquerait que la séparation observée entre les
courbes peut être décrite comme un décalage horizontal entre deux exponentielles
de pente commune.

Une proximité entre
\(\Delta t_{ij}^{\mathrm{fit}}\) et
\(\Delta t_{ij}^{F}\) indiquerait que la différence entre les amplitudes initiales
extrapolées est cohérente avec les niveaux de perturbation estimés à partir de
la force.

Si les trois valeurs sont proches, l’effet de la taille de dexel peut être
interprété, dans cette première hypothèse, comme une modification du niveau
initial de perturbation sans changement significatif du taux de croissance.

Dans le cas contraire, l’hypothèse
\(\beta_i^{\mathrm{fit}}=\beta_j^{\mathrm{fit}}\) devra être abandonnée et le cas
général de taux différents devra être étudié.

\begin{figure}[H]
    \centering
    \fbox{
        \parbox[c][5cm][c]{0.85\textwidth}{
            \centering
            \textbf{Figure à ajouter}\\[0.3cm]
            Comparaison de
            \(\Delta t_{ij}^{\mathrm{mes}}\),
            \(\Delta t_{ij}^{\mathrm{fit}}\)
            et
            \(\Delta t_{ij}^{F}\)
            sous l’hypothèse d’un taux de croissance commun.
        }
    }
    \caption{Évaluation de l’hypothèse de taux de croissance identiques.}
    \label{fig:evaluation-beta-identiques}
\end{figure}



%==============================================
% == Section 7 ==========================================
%==============================================

\section{Cas général : taux de croissance différents}
\label{sec:beta-differents}

Les résultats précédents reposaient sur l’hypothèse

\begin{equation}
    \beta_i^{\mathrm{fit}}
    =
    \beta_j^{\mathrm{fit}}.
\end{equation}

Cette hypothèse doit être abandonnée lorsque les pentes obtenues dans les
intervalles d’ajustement présentent une différence non négligeable. Le décalage
temporel ne dépend alors plus uniquement des amplitudes initiales extrapolées.
Il dépend également de la différence entre les taux de croissance.

Dans cette section, aucune égalité n’est donc imposée entre

\begin{equation}
    \beta_i^{\mathrm{fit}}
    \qquad \text{et} \qquad
    \beta_j^{\mathrm{fit}}.
\end{equation}

\subsection{Décomposition du décalage obtenu par ajustement}
\label{subsec:decomposition-fit-beta-differents}

Le décalage général obtenu à partir des ajustements a été établi à la
\ref{eq:decalage-fit-general} :

\begin{equation}
    \Delta t_{ij}^{\mathrm{fit}}
    =
    \frac{
        \ln A_{\mathrm{ref}}
        -
        \ln P_j^{\mathrm{fit}}
    }{
        \beta_j^{\mathrm{fit}}
    }
    -
    \frac{
        \ln A_{\mathrm{ref}}
        -
        \ln P_i^{\mathrm{fit}}
    }{
        \beta_i^{\mathrm{fit}}
    }.
    \label{eq:dt-fit-general-rappel}
\end{equation}

Afin de séparer l’effet des amplitudes initiales de celui des taux de
croissance, le terme

\begin{equation}
    \frac{
        \ln A_{\mathrm{ref}}
        -
        \ln P_i^{\mathrm{fit}}
    }{
        \beta_j^{\mathrm{fit}}
    }
\end{equation}

est ajouté puis soustrait. Cette opération ne modifie pas la valeur du
décalage, mais permet de regrouper les termes selon leur origine :

\begin{align}
    \Delta t_{ij}^{\mathrm{fit}}
    &=
    \frac{
        \ln A_{\mathrm{ref}}
        -
        \ln P_j^{\mathrm{fit}}
    }{
        \beta_j^{\mathrm{fit}}
    }
    -
    \frac{
        \ln A_{\mathrm{ref}}
        -
        \ln P_i^{\mathrm{fit}}
    }{
        \beta_j^{\mathrm{fit}}
    }
    +
    \frac{
        \ln A_{\mathrm{ref}}
        -
        \ln P_i^{\mathrm{fit}}
    }{
        \beta_j^{\mathrm{fit}}
    }
    -
    \frac{
        \ln A_{\mathrm{ref}}
        -
        \ln P_i^{\mathrm{fit}}
    }{
        \beta_i^{\mathrm{fit}}}.
    \label{eq:dt-fit-ajout-soustraction}
\end{align}

En regroupant les deux premières fractions, on obtient

\begin{align}
    \Delta t_{ij}^{\mathrm{fit}}
    &=
    \frac{
        \ln P_i^{\mathrm{fit}}
        -
        \ln P_j^{\mathrm{fit}}
    }{
        \beta_j^{\mathrm{fit}}
    }
    \nonumber\\
    &\quad+
    \left(
        \ln A_{\mathrm{ref}}
        -
        \ln P_i^{\mathrm{fit}}
    \right)
    \left(
        \frac{1}{\beta_j^{\mathrm{fit}}}
        -
        \frac{1}{\beta_i^{\mathrm{fit}}}
    \right).
    \label{eq:decomposition-dt-fit}
\end{align}

Le décalage peut ainsi être écrit comme la somme de deux contributions :

\begin{equation}
    \Delta t_{ij}^{\mathrm{fit}}
    =
    \Delta t_{P,ij}^{\mathrm{fit}}
    +
    \Delta t_{\beta,ij}^{\mathrm{fit}}.
    \label{eq:somme-contributions-fit}
\end{equation}

La contribution associée aux amplitudes initiales est définie par

\begin{equation}
    \Delta t_{P,ij}^{\mathrm{fit}}
    =
    \frac{
        \ln P_i^{\mathrm{fit}}
        -
        \ln P_j^{\mathrm{fit}}
    }{
        \beta_j^{\mathrm{fit}}
    }
    =
    \frac{1}{\beta_j^{\mathrm{fit}}}
    \ln\left(
        \frac{P_i^{\mathrm{fit}}}
        {P_j^{\mathrm{fit}}}
    \right).
    \label{eq:contribution-P-fit}
\end{equation}

Cette contribution représente le décalage qui serait produit par la différence
entre les amplitudes initiales extrapolées, en utilisant
\(\beta_j^{\mathrm{fit}}\) comme taux de référence.

La contribution associée aux taux de croissance est définie par

\begin{equation}
    \Delta t_{\beta,ij}^{\mathrm{fit}}
    =
    \left(
        \ln A_{\mathrm{ref}}
        -
        \ln P_i^{\mathrm{fit}}
    \right)
    \left(
        \frac{1}{\beta_j^{\mathrm{fit}}}
        -
        \frac{1}{\beta_i^{\mathrm{fit}}}
    \right).
    \label{eq:contribution-beta-fit}
\end{equation}

Le premier facteur,

\begin{equation}
    \ln A_{\mathrm{ref}}
    -
    \ln P_i^{\mathrm{fit}}
    =
    \ln\left(
        \frac{A_{\mathrm{ref}}}
        {P_i^{\mathrm{fit}}}
    \right),
\end{equation}

représente la distance logarithmique entre l’amplitude initiale extrapolée du
cas \(i\) et le niveau de référence. Le second facteur représente la différence
entre les inverses des taux de croissance.

Lorsque

\begin{equation}
    \beta_i^{\mathrm{fit}}
    =
    \beta_j^{\mathrm{fit}},
\end{equation}

la contribution \(\Delta t_{\beta,ij}^{\mathrm{fit}}\) devient nulle et la
formulation du cas de taux identiques est retrouvée.

Il faut noter que cette décomposition utilise le cas \(i\) comme référence
algébrique et \(\beta_j^{\mathrm{fit}}\) comme taux de référence pour la
contribution d’amplitude. Une autre décomposition pourrait être obtenue en
inversant les rôles de \(i\) et \(j\). La valeur totale
\(\Delta t_{ij}^{\mathrm{fit}}\) reste cependant inchangée.

\subsection{Décomposition du décalage reconstruit à partir de la force}
\label{subsec:decomposition-force-beta-differents}

Le décalage général reconstruit à partir de la force a été défini à la
\ref{eq:decalage-force-general}. Comme pour le décalage issu des ajustements,
il peut être séparé en une contribution associée au niveau initial de
perturbation et une contribution associée à la différence entre les taux de
croissance.

En utilisant l’ordonnée à l’origine reconstruite à partir de la force,
\(c_i^{F}\), le décalage général possède la même structure que celui obtenu par
ajustement :

\begin{equation}
    \Delta t_{ij}^{F}
    =
    \frac{
        \ln A_{\mathrm{ref}}-c_j^{F}
    }{
        \beta_j^{\mathrm{fit}}
    }
    -
    \frac{
        \ln A_{\mathrm{ref}}-c_i^{F}
    }{
        \beta_i^{\mathrm{fit}}
    }.
    \label{eq:dt-force-forme-ciF}
\end{equation}

En ajoutant puis en soustrayant

\begin{equation}
    \frac{
        \ln A_{\mathrm{ref}}-c_i^{F}
    }{
        \beta_j^{\mathrm{fit}}
    },
\end{equation}

on obtient

\begin{align}
    \Delta t_{ij}^{F}
    &=
    \frac{
        c_i^{F}-c_j^{F}
    }{
        \beta_j^{\mathrm{fit}}
    }
    \nonumber\\
    &\quad+
    \left(
        \ln A_{\mathrm{ref}}-c_i^{F}
    \right)
    \left(
        \frac{1}{\beta_j^{\mathrm{fit}}}
        -
        \frac{1}{\beta_i^{\mathrm{fit}}}
    \right).
    \label{eq:decomposition-dt-force-ciF}
\end{align}

Le décalage reconstruit à partir de la force est donc écrit sous la forme

\begin{equation}
    \Delta t_{ij}^{F}
    =
    \Delta t_{P,ij}^{F}
    +
    \Delta t_{\beta,ij}^{F}.
    \label{eq:somme-contributions-force}
\end{equation}

La première contribution est

\begin{equation}
    \Delta t_{P,ij}^{F}
    =
    \frac{
        c_i^{F}-c_j^{F}
    }{
        \beta_j^{\mathrm{fit}}
    }.
    \label{eq:contribution-P-force-ciF}
\end{equation}

D’après l’expression de \(c_i^{F}\) établie dans la section précédente,

\begin{equation}
    c_i^{F}
    =
    \ln P_{h,i}^{F}
    -
    \ln K_f
    -
    \ln C_i
    -
    \beta_i^{\mathrm{fit}}t_{s,i},
    \label{eq:ciF-rappel-section7}
\end{equation}

et, de manière analogue,

\begin{equation}
    c_j^{F}
    =
    \ln P_{h,j}^{F}
    -
    \ln K_f
    -
    \ln C_j
    -
    \beta_j^{\mathrm{fit}}t_{s,j}.
    \label{eq:cjF-rappel-section7}
\end{equation}

En substituant ces expressions dans la
\ref{eq:contribution-P-force-ciF}, le coefficient commun \(K_f\) disparaît :

\begin{align}
    \Delta t_{P,ij}^{F}
    &=
    \frac{1}{\beta_j^{\mathrm{fit}}}
    \left[
        \ln P_{h,i}^{F}
        -
        \ln P_{h,j}^{F}
        +
        \ln C_j
        -
        \ln C_i
    \right]
    -
    \frac{\beta_i^{\mathrm{fit}}}
    {\beta_j^{\mathrm{fit}}}
    t_{s,i}
    +
    t_{s,j},
    \nonumber\\
    &=
    \frac{1}{\beta_j^{\mathrm{fit}}}
    \ln\left(
        \frac{
            P_{h,i}^{F}C_j
        }{
            P_{h,j}^{F}C_i
        }
    \right)
    -
    \frac{\beta_i^{\mathrm{fit}}}
    {\beta_j^{\mathrm{fit}}}
    t_{s,i}
    +
    t_{s,j}.
    \label{eq:contribution-P-force}
\end{align}

Cette contribution contient trois effets :

\begin{description}
    \item[Rapport des proxies]
    la différence entre les niveaux \(P_{h,i}^{F}\) et \(P_{h,j}^{F}\) mesurés
    à partir de la force normalisée ;

    \item[Facteurs \(C_i\) et \(C_j\)]
    la différence éventuelle entre les relations de proportionnalité associées
    aux deux cas ;

    \item[Correction temporelle]
    la prise en compte des instants \(t_{s,i}\) et \(t_{s,j}\) auxquels les
    proxies ont été évalués.
\end{description}

La seconde contribution est

\begin{equation}
    \Delta t_{\beta,ij}^{F}
    =
    \left(
        \ln A_{\mathrm{ref}}-c_i^{F}
    \right)
    \left(
        \frac{1}{\beta_j^{\mathrm{fit}}}
        -
        \frac{1}{\beta_i^{\mathrm{fit}}}
    \right).
    \label{eq:contribution-beta-force-ciF}
\end{equation}

En remplaçant \(c_i^{F}\) par son expression développée, il vient

\begin{align}
    \Delta t_{\beta,ij}^{F}
    &=
    \Bigl(
        \ln A_{\mathrm{ref}}
        -
        \ln P_{h,i}^{F}
        +
        \ln K_f
        +
        \ln C_i
        +
        \beta_i^{\mathrm{fit}}t_{s,i}
    \Bigr)
    \nonumber\\
    &\quad\times
    \left(
        \frac{1}{\beta_j^{\mathrm{fit}}}
        -
        \frac{1}{\beta_i^{\mathrm{fit}}}
    \right).
    \label{eq:contribution-beta-force}
\end{align}

Cette contribution quantifie la partie du décalage reconstruit qui apparaît
uniquement parce que les deux taux de croissance sont différents.

Si le facteur de proportionnalité peut être considéré comme commun,

\begin{equation}
    C_i\approx C_j,
\end{equation}

la contribution associée au niveau initial se simplifie en

\begin{equation}
    \Delta t_{P,ij}^{F}
    \approx
    \frac{1}{\beta_j^{\mathrm{fit}}}
    \ln\left(
        \frac{P_{h,i}^{F}}
        {P_{h,j}^{F}}
    \right)
    -
    \frac{\beta_i^{\mathrm{fit}}}
    {\beta_j^{\mathrm{fit}}}
    t_{s,i}
    +
    t_{s,j}.
    \label{eq:contribution-P-force-simplifiee}
\end{equation}

\subsection{Comparaison des contributions}
\label{subsec:comparaison-contributions-beta-differents}

La décomposition précédente permet de distinguer deux mécanismes susceptibles
de produire le décalage observé :

\begin{description}
    \item[Effet du niveau initial]
    représenté par
    \(\Delta t_{P,ij}^{\mathrm{fit}}\) et
    \(\Delta t_{P,ij}^{F}\) ;

    \item[Effet du taux de croissance]
    représenté par
    \(\Delta t_{\beta,ij}^{\mathrm{fit}}\) et
    \(\Delta t_{\beta,ij}^{F}\).
\end{description}

Pour les ajustements, la cohérence de la décomposition doit être vérifiée par

\begin{equation}
    \Delta t_{ij}^{\mathrm{fit}}
    =
    \Delta t_{P,ij}^{\mathrm{fit}}
    +
    \Delta t_{\beta,ij}^{\mathrm{fit}}.
    \label{eq:verification-decomposition-fit}
\end{equation}

Pour la reconstruction fondée sur la force, elle doit être vérifiée par

\begin{equation}
    \Delta t_{ij}^{F}
    =
    \Delta t_{P,ij}^{F}
    +
    \Delta t_{\beta,ij}^{F}.
    \label{eq:verification-decomposition-force}
\end{equation}

Les valeurs totales sont ensuite comparées au décalage directement mesuré
\(\Delta t_{ij}^{\mathrm{mes}}\).

Si

\begin{equation}
    \left|
        \Delta t_{P,ij}^{\mathrm{fit}}
    \right|
    \gg
    \left|
        \Delta t_{\beta,ij}^{\mathrm{fit}}
    \right|,
\end{equation}

le décalage est principalement associé à une différence entre les amplitudes
initiales extrapolées.

À l’inverse, si

\begin{equation}
    \left|
        \Delta t_{\beta,ij}^{\mathrm{fit}}
    \right|
    \not\ll
    \left|
        \Delta t_{P,ij}^{\mathrm{fit}}
    \right|,
\end{equation}

la différence entre les taux de croissance contribue de manière significative
au décalage et ne peut pas être négligée.

La comparaison entre les contributions issues des ajustements et celles
reconstruites à partir de la force permet finalement d’évaluer si la différence
entre les niveaux initiaux est cohérente avec la perturbation estimée à partir
de la force normalisée.

\begin{figure}[H]
    \centering
    \fbox{
        \parbox[c][5cm][c]{0.85\textwidth}{
            \centering
            \textbf{Figure à ajouter}\\[0.3cm]
            Comparaison de
            \(\Delta t_{P,ij}^{\mathrm{fit}}\),
            \(\Delta t_{\beta,ij}^{\mathrm{fit}}\),
            \(\Delta t_{P,ij}^{F}\),
            \(\Delta t_{\beta,ij}^{F}\)
            et \(\Delta t_{ij}^{\mathrm{mes}}\)
            pour chaque paire de tailles de dexel.
        }
    }
    \caption{Décomposition du décalage en contributions associées au niveau initial et au taux de croissance.}
    \label{fig:decomposition-decalage-beta-differents}
\end{figure}






\end{document}